% Part: many-valued-logic
% Chapter: infinite-valued-logics
% Section: goedel

\documentclass[../../../include/open-logic-section]{subfiles}

\begin{document}

\olfileid{mvl}{inf}{god}

\olsection{G\"odel যুক্তিবিদ্যাসমূহ}

\begin{editorial}
  এটি অসীমমানী G\"odel যুক্তিবিদ্যার একটি সংক্ষিপ্ত প্রাথমিক খসড়া।
\end{editorial}

\begin{defn}\ollabel{def:goedel} অসীমমানী G\"odel
যুক্তিবিদ্যা~$\LogGod[\infty]$ নিচের ম্যাট্রিক্স দিয়ে সংজ্ঞায়িত:
\begin{enumerate}
  \item $\lfalse$, $\lnot$, $\land$, $\lor$, $\lif$-সহ প্রমিত
  বচনগত ভাষা $\Lang L_0$।
  \item সত্যমানের সেট $V_\infty$।
  \item $1$ একমাত্র মনোনীত মান; অর্থাৎ, $V^+ = \{1\}$।
  \item সত্যমান-অপেক্ষকগুলি নিচের অপেক্ষকগুলির মাধ্যমে নির্ধারিত:
  \begin{align*}
    \tf{\lfalse} & = 0\\
    \tf{\lnot}[\LogGod](x) & = \begin{cases}
      1 & \text{if } x =0\\
      0 & \text{otherwise}
    \end{cases}\\
    \tf{\land}[\LogGod](x,y) & = \min(x,y)\\
    \tf{\lor}[\LogGod](x,y) & = \max(x,y)\\
    \tf{\lif}[\LogGod](x,y) & = \begin{cases}
      1 & \text{if } x \le y\\
      y & \text{otherwise.}
    \end{cases}
    \end{align*}
\end{enumerate}
$m$-মানী G\"odel যুক্তিবিদ্যাও একইভাবে সংজ্ঞায়িত, শুধু
$V = V_m$ নেওয়া হয়।
\end{defn}
% BN-SRC-326 TARGET-CORRECTION-BEGIN
\par\smallskip\noindent\textnormal{\small\textbf{উৎস-সংশোধন (BN-SRC-326):} হিমায়িত উৎসের \texttt{cases} পরিবেশে নঞর্থকরণের ফল হিসেবে \texttt{\$1\$} ও \texttt{\$0\$} লেখা আছে। ইতিমধ্যে গণিত-মোডে থাকা এই পরিবেশে ভেতরের ডলারচিহ্ন TeX-এর মোড ভাঙে; বাংলা লক্ষ্যে কেবল সেগুলি বাদ দিয়ে যথাক্রমে $1$ ও~$0$ রাখা হয়েছে।} % BN-SRC-326 TARGET-CORRECTION-END

\begin{prop}
  \olref[thr][god]{defn:goedel} অনুযায়ী সংজ্ঞায়িত
  $\LogGod[3]$ যুক্তিবিদ্যা এবং \olref{def:goedel} অনুযায়ী
  সংজ্ঞায়িত $\LogGod[3]$ যুক্তিবিদ্যা একই।
\end{prop}

\begin{proof}
  \olref[thr][god]{defn:goedel}-এ দেওয়া সংযোজকগুলির সত্যসারণি
  \olref{def:goedel}-এর সমীকরণ থেকে পাওয়া নিচের সত্যসারণিগুলির
  সঙ্গে তুলনা করলেই তা দেখা যায়:
  \begin{center}
    \begin{tabular}{c|c}
      $\tf{\lnot}[\LogGod[3]]$ & \\
      \hline
      $1$ & $0$ \\
      $1/2$ & $0$ \\
      $0$ & $1$
    \end{tabular}
    \quad
    \begin{tabular}{c|ccc}
      $\tf{\land}[\LogGod]$ & $1$ & $1/2$ & $0$ \\
      \hline
      $1$ & $1$ & $1/2$ & $0$ \\
      $1/2$ & $1/2$ & $1/2$ & $0$\\
      $0$ & $0$ & $0$ & $0$
    \end{tabular}
    \\[2ex]
    \begin{tabular}{c|ccc}
      $\tf{\lor}[\LogGod]$ & $1$ & $1/2$ & $0$ \\
      \hline
      $1$ & $1$ & $1$ & $1$ \\
      $1/2$ & $1$ & $1/2$ & $1/2$ \\
      $0$ & $1$ & $1/2$ & $0$
    \end{tabular}
    \quad
    \begin{tabular}{c|ccc}
      $\tf{\lif}[\LogGod]$ & $1$ & $1/2$ & $0$ \\
      \hline
      $1$ & $1$ & $1/2$ & $0$ \\
      $1/2$ & $1$ & $1$ & $0$ \\
      $0$ & $1$ & $1$ & $1$
    \end{tabular}
  \end{center}
\end{proof}

\begin{prop}\ollabel{prop:god-infty-m}
  যদি $\Gamma \Entails[\LogGod[\infty]] !B$ হয়, তবে সব~$m \ge 2$-এর
  জন্য $\Gamma \Entails[\LogGod[m]] !B$।
\end{prop}

\begin{proof}
  অনুশীলনী।
\end{proof}

\begin{prob}
  \olref[mvl][inf][god]{prop:god-infty-m} প্রমাণ করো।
\end{prob}

সসীম $\Gamma$-এর ক্ষেত্রে এর বিপরীত দাবিটিও সত্য।
% BN-SRC-328 TARGET-CORRECTION-BEGIN
\par\smallskip\noindent\textnormal{\small\textbf{উৎস-সংশোধন (BN-SRC-328):} হিমায়িত উৎসে $\Gamma$ সসীম হওয়ার শর্ত নেই, অথচ অসীম পূর্বধারণাসেটে দাবিটি ব্যর্থ। $n\ge1$-এর জন্য $\Gamma$-তে $p_n\lif q$ এবং $((p_{n+1}\lif p_n)\lif p_n)$ রাখি। কোনো সসীম $V_m$-এ $q<1$ ধরে সব পূর্বধারণা সত্য হলে প্রথমগুলি থেকে $p_n\le q<1$ এবং দ্বিতীয়গুলি থেকে $p_n<p_{n+1}$ হয়; সসীম $V_m$-এ এমন অসীম ক্রম অসম্ভব। ফলে সব $m$-এ $\Gamma\Entails[\LogGod[m]]q$। কিন্তু $V_\infty$-এ $q=1/2$ ও $p_n=n/(2(n+1))$ দিলে সব পূর্বধারণা সত্য, $q$ মনোনীত নয়। সসীম $\Gamma$-এর কোনো প্রতিবিপরীত মূল্যায়নে যতগুলি মূলদ মান ব্যবহৃত হয়, সবগুলি একটি $V_m$-এ থাকে; তাই সসীম ক্ষেত্রে বিপরীত দাবিটি সত্য।} % BN-SRC-328 TARGET-CORRECTION-END

$\LogGod[3]$-এর মতো $\LogGod[\infty]$-এও সব স্বজ্ঞাবাদীভাবে সিদ্ধ
!!{formula}s সর্বতঃসত্য, এবং আগের অ-সর্বতঃসত্য সূত্রগুলিও
$\LogGod[\infty]$-এ অ-সর্বতঃসত্য:
\begin{align*}
  & p \lor \lnot p && (p \lif q) \lif (\lnot p \lor q) \\
  & \lnot\lnot p \lif p && \lnot(\lnot p \land \lnot q) \lif (p \lor q) \\
  & ((p \lif q) \lif p) \lif p && \lnot(p \lif q) \lif (p \land \lnot q)
\end{align*}
স্বজ্ঞাবাদীভাবে সিদ্ধ নয়, অথচ $\LogGod[3]$-এ সর্বতঃসত্য---এমন
!!{formula}-এর উদাহরণ $(p \lif q) \lor (q \lif p)$-ও
$\LogGod[\infty]$-এ সর্বতঃসত্য। আরও স্পষ্টভাবে,
$\LogGod[\infty]$-কে স্বজ্ঞাবাদী যুক্তিবিদ্যার সঙ্গে
$(!A \lif !B) \lor (!B \lif !A)$ প্রকল্পটি যোগ করে পাওয়া
যুক্তিবিদ্যা হিসেবে চিহ্নিত করা যায়। মাইকেল ডামেট এটি দেখিয়েছিলেন;
তাই $\LogGod[\infty]$-কে প্রায়ই G\"odel--Dummett
যুক্তিবিদ্যা~$\Log{LC}$ বলা হয়।

\begin{prob}
  দেখাও যে $(p \lif q) \lor (q \lif p)$ সূত্রটি
  $\LogGod[\infty]$-এ সর্বতঃসত্য।
\end{prob}

\begin{prob}
  দেখাও যে $(p \lif q) \lor (q \lif r) \lor (r \lif s)$ সূত্রটি
  $\LogGod[3]$-এ সর্বতঃসত্য হলেও $\LogGod[\infty]$-এ
  সর্বতঃসত্য নয়।
\end{prob}

\end{document}
